% Modelo de slides para projetos de disciplinas do Abel
\documentclass[10pt]{beamer}

\usetheme[progressbar=frametitle]{metropolis}
\usepackage{appendixnumberbeamer}
\usepackage[]{natbib}
\usepackage{todonotes}
\presetkeys{todonotes}{inline}{}


% \usepackage[nottoc]{tocbibind}
% \usepackage{biblatex}
\renewcommand\appendixname{Appendix}

\usepackage{subfig}
\usepackage{tikz}
\usepackage{pgfplots}
\usepgfplotslibrary{fillbetween}
% \pgfplotsset{compat=1.16}
\allowdisplaybreaks

\usepackage{xcolor}

\definecolor{ultramarine}{RGB}{35,174,67}
\definecolor{f3}{HTML}{61D77D}

\makeatletter
\def\mathcolor#1#{\@mathcolor{#1}}
\def\@mathcolor#1#2#3{%
  \protect\leavevmode
  \begingroup
    \color#1{#2}#3%
  \endgroup
}
\makeatother

\usepackage[usenames,dvipsnames]{pstricks}

\usepackage{algorithm}
\usepackage{algorithmic}

\usepackage[usenames,dvipsnames]{pstricks}
\usepackage{epsfig}
\usepackage{pst-grad} % For gradients
\usepackage{pst-plot} % For axes
\usepackage[space]{grffile} % For spaces in paths
\usepackage{etoolbox} % For spaces in paths
\makeatletter % For spaces in paths
\patchcmd\Gread@eps{\@inputcheck#1 }{\@inputcheck"#1"\relax}{}{}
\makeatother
% User Packages:
\usepackage{amssymb,amsthm,amsmath}
\usepackage{mathtext}

\usepackage{epsfig}
\usepackage{pst-grad} % For gradients
\usepackage{pst-plot} % For axes
\usepackage[space]{grffile} % For spaces in paths
\usepackage{etoolbox} % For spaces in paths
\makeatletter % For spaces in paths
\patchcmd\Gread@eps{\@inputcheck#1 }{\@inputcheck"#1"\relax}{}{}
\makeatother
\usepackage{booktabs}
\usepackage[scale=2]{ccicons}

\usepackage{xspace}
\newcommand{\themename}{\textbf{\textsc{metropolis}}\xspace}

\DeclareMathOperator*{\arginf}{arg\,inf}

\makeatletter
\setbeamertemplate{title page}{
  \begin{minipage}[b][\paperheight]{\textwidth}
    \centering  % <-- Center here
    \ifx\inserttitlegraphic\@empty\else\usebeamertemplate*{title graphic}\fi
    \vfill%
    \ifx\inserttitle\@empty\else\usebeamertemplate*{title}\fi
    \ifx\insertsubtitle\@empty\else\usebeamertemplate*{subtitle}\fi
    \usebeamertemplate*{title separator}
    \ifx\beamer@shortauthor\@empty\else\usebeamertemplate*{author}\fi
    \ifx\insertdate\@empty\else\usebeamertemplate*{date}\fi
    \ifx\insertinstitute\@empty\else\usebeamertemplate*{institute}\fi
    \vfill
    \vspace*{1mm}
  \end{minipage}
}

\title{Online Gradient Descent for Online Portfolio Optimization with Transaction Costs}
% \subtitle{Subtítulo}
% \date{\today}
\date{}
\author{\bf{ Martino Nicol\a'o Bernasconi de Luca }\\ \footnotesize Supervisor: Francesco Trov\a'o \\
Co-supervisors: Marcello Restelli, Edoardo Vittori}

\institute[Politecnico di Milano]{\includegraphics[width=3cm]{./img/Logo_Politecnico_Milano.png}}
% \titlegraphic{\hfill\includegraphics[height=1.5cm]{logo.pdf}}

\begin{document}

\begin{frame} 
\centering 
\titlepage 
\end{frame}

\begin{frame}{Table of contents}
  \setbeamertemplate{section in toc}[sections numbered]
  \tableofcontents[hideallsubsections]
\end{frame}

\section{Context}

\begin{frame}{Online Portfolio Optimization}


\begin{columns}[T] % align columns
\begin{column}{.41\textwidth}
\onslide<1->{
\vspace{-2em}
\begin{figure}
\centering
 \scalebox{.43}{\hspace{1.5em}\begin{pspicture}(0,-6.6)(10.4,6.6)
\definecolor{colour0}{rgb}{0.14,0.69,0.27}
\rput(5.0,-0.6){\includegraphics{/Users/Martino/Documents/GitHub/thesis/presentation/img/stocks.png}}
\psline[linecolor=black, linewidth=0.04](0.0,0.57)(10.4,0.57)
\psframe[linecolor=white, linewidth=0.002, fillstyle=solid,fillcolor=colour0, dimen=outer](9.2,2.2)(8.4,0.6)
\rput[bl](0.8,4){$20\%$}
\rput[bl](3.2,4){$15\%$}
\rput[bl](6.0,4){$40\%$}
\rput[bl](8.8,4){$25\%$}
\psframe[linecolor=white, linewidth=0.002, fillstyle=solid,fillcolor=colour0, dimen=outer](6.8,3.4)(6.0,0.6)
\psframe[linecolor=white, linewidth=0.002, fillstyle=solid,fillcolor=colour0, dimen=outer](4.0,1.4)(3.2,0.6)
\psframe[linecolor=white, linewidth=0.002, fillstyle=solid,fillcolor=colour0, dimen=outer](1.6,1.8)(0.8,0.6)
%\psframe[linecolor=black, linewidth=0.04, dimen=outer](7.6,6.6)(2.4,5.4)
%\rput[bl](4.4,5.8){Agent: $\mathcal A$}
%\psline[linecolor=black, linewidth=0.04, arrowsize=0.05291667cm 2.0,arrowlength=1.4,arrowinset=0.0]{->}(4.8,5.4)(4.8,4.2)
%\psframe[linecolor=black, linewidth=0.04, dimen=outer](7.2,-5.4)(2.4,-6.6)
%\rput[bl](4.0,-6.2){Adversary}
%\psline[linecolor=black, linewidth=0.04, arrowsize=0.05291667cm 2.0,arrowlength=1.4,arrowinset=0.0]{->}(4.8056273,-5.3999734)(4.7943726,-4.2000265)
\end{pspicture}}
\end{figure}
}
\vspace{-6em}
\begin{figure}[t!]
    \scalebox{.43}{\begin{tikzpicture}

\colorlet{colorUP}{red}
\colorlet{colorOGD}{green!50.0!black}
\colorlet{colorONS}{blue}
\definecolor{colorOLU}{rgb}{0.75,0.75,0}
\pgfkeys{/pgf/number format/.cd,1000 sep={}}
\pgfplotsset{every tick label/.append style={font=\tiny}}
\begin{axis}[
height = 8cm,
width = 10cm,
legend cell align={left},
legend style={at={(0.03,0.97)}, anchor=north west, draw=white!80.0!black},
tick align=outside,
tick pos=left,
x grid style={white!69.01960784313725!black},
xlabel={$time$},
xmin=0.5, xmax=5700.5,
ymin=0, ymax=45,
xtick style={color=black},
xtick = {1, 1000, 2000, ..., 5000, 6000},
y grid style={white!69.01960784313725!black},
ylabel={$Wealth$},
ytick style={color=black},
yticklabel style={rotate=90},
]
\addplot [thick, colorOGD, 
%mark = o,
draw opacity=0.4,
mark repeat=20,mark phase=7]
table {./figures/fig_tmp/data_B.txt};
\addlegendentry{\small $S_1$}
\addplot [thick, colorUP, %mark = square,
draw opacity=0.4,
mark repeat=20,mark phase=7]
table {./figures/fig_tmp/data_C.txt};
\addlegendentry{\small $S_2$}
\addplot [semithick, colorOLU, %mark = diamond,
draw opacity=0.4,
mark repeat=20,mark phase=7]
table {./figures/fig_tmp/data_H.txt};
\addlegendentry{\small $S_3$}
\addplot [semithick, colorONS, %mark = triangle,
draw opacity=0.4,
mark repeat=20,mark phase=7]
table {./figures/fig_tmp/data_I.txt};
\addlegendentry{\small $S_4$}
\addplot [thick, black, %mark = +,
mark repeat=20,mark phase=7]
table {./figures/fig_tmp/data_P.txt};
\addlegendentry{Portfolio}
\end{axis}

\end{tikzpicture}
}
\end{figure}
\end{column}
\begin{column}{.45\textwidth}


\onslide<2->{
\vspace{2em}
\begin{itemize}
  \item Modern Portfolio Optimization:
  \begin{itemize}
    \item Statistical assumptions on the stock dynamics
    \item Backward looking
  \end{itemize}

  \item Online Portfolio Optimization:
  \begin{itemize}
    \item Adversarial (no assumptions on the distribution)
    \item From Online Learning field
  \end{itemize}
\end{itemize}
}

\end{column}

\end{columns}
\end{frame}



\begin{frame}[fragile]{Context: Online Learning with Experts Advice}

\begin{figure}
\centering
 \scalebox{.45}{\begin{pspicture}(0,-3.42)(8.18,3.42)
\psframe[linecolor=black, linewidth=0.04, dimen=outer, framearc=0.1](6.54,0.9)(2.78,-0.36)
\psframe[linecolor=black, linewidth=0.04, dimen=outer, framearc=0.1](6.52,3.42)(2.82,2.3)
\rput[bl](3.76,2.7){Agent: $\mathcal A$}
\rput[bl](3.68,0.12){Environment}
\psline[linecolor=black, linewidth=0.036, arrowsize=0.05291667cm 2.0,arrowlength=1.4,arrowinset=0.0]{->}(6.54,2.8)(7.66,2.8)(7.66,0.2)(6.54,0.22)
\rput[bl](7.9,1.44){$x_t$}
\psline[linecolor=black, linewidth=0.036, arrowsize=0.05291667cm 2.0,arrowlength=1.4,arrowinset=0.0]{->}(2.78,0.3)(1.46,0.3)(1.46,2.88)(2.8267796,2.8677726)
\rput[bl](0.0,1.2){$f(x_t,y_t)$}
\psframe[linecolor=black, linewidth=0.036, dimen=outer, framearc=0.1](6.2,-2.2)(3.16,-3.42)
\rput[bl](3.78,-2.92){Experts: $\mathcal E$}
\rput[bl](2.96,-1.36){$x_{e,t}$}
\rput[bl](5.26,-1.34){$f(x_{e,t},y_t)$}
\psline[linecolor=black, linewidth=0.036, arrowsize=0.05291667cm 2.0,arrowlength=1.4,arrowinset=0.0]{->}(3.8,-2.22)(3.8,-0.36)
\psline[linecolor=black, linewidth=0.036, arrowsize=0.05291667cm 2.0,arrowlength=1.4,arrowinset=0.0]{->}(5.18,-0.4)(5.18,-2.22)
\end{pspicture}
}
\end{figure}

\onslide<5->{
\begin{itemize}
% \item $L_T(\mathcal A):=\sum\limits_{t=1}^Tf_t(x_t)$
% \item $L_{T}^*:=\inf\limits_{e\in \mathcal E}\sum\limits_{t=1}^Tf_t(x_{e,t})$
\item Regret: $R_T=\sum\limits_{t=1}^Tf_t(x_t)-\inf\limits_{e\in \mathcal E}\sum\limits_{t=1}^Tf_t(x_{e,t})$
\begin{itemize}
  \item No regret: $R_T=o(T)$ for any sequence $f_1,f_2,\ldots$
\end{itemize}
\item Per-round Computational Complexity
\end{itemize}
}
\onslide<6->{
\begin{figure}
\centering
 \scalebox{.4}{\begin{pspicture}(0,-1.39)(8.72,1.39)
\definecolor{colour0}{rgb}{0.38039216,0.84313726,0.49019608}
\definecolor{colour1}{rgb}{0.38039216,0.627451,0.84313726}
\definecolor{colour2}{rgb}{0.87058824,0.47058824,0.12156863}
\pstriangle[linecolor=black, linewidth=0.04, fillstyle=solid,fillcolor=colour0, opacity=0.69803923, dimen=outer](4.34,-1.39)(1.2908762,1.3750794)
\psline[linecolor=black, linewidth=0.04](0.0,0.06460317)(8.72,0.06460317)
\psframe[linecolor=black, linewidth=0.04, fillstyle=solid,fillcolor=colour1, opacity=0.6784314, dimen=outer](2.6344411,1.39)(0.0,0.03809524)
\psframe[linecolor=black, linewidth=0.04, fillstyle=solid,fillcolor=colour2, opacity=0.73333335, dimen=outer](8.72,1.39)(5.822115,0.03809524)
\rput[bl](0.28,0.85){\small Computational}
\rput[bl](0.3,0.41){Complexity}
\rput[bl](6.68,0.59){Regret}
\end{pspicture}}
\end{figure}
}

% \begin{columns}[T] % align columns
% \begin{column}{.41\textwidth}
% % \color{red}\rule{\linewidth}{4pt}

% \end{column}%
% \hfill%
% \begin{column}{.55\textwidth}
% % \color{blue}\rule{\linewidth}{4pt}


% \end{column}%
% \end{columns}

% \onslide<6->{
%         \begin{center}
%             \begin{minipage}{10cm}
%             \metroset{block=fill}
%                 \begin{exampleblock}{No Regret}
%                 % \centering
%                 $\mathcal A$ learns if has per round zero regret asymptotically: $R_T=o(T)$ for any sequence $f_1,f_2,\ldots,$ of losses
%                 \end{exampleblock}
%             \end{minipage}
%         \end{center}}

\end{frame}

\section{Problem Formulation}

\begin{frame}{From Online Learning to Portfolio}

\begin{itemize}
\item $\mathbf x_t\in\Delta_{N-1}$ is the portfolio allocation
\item $\mathbf y_t=\left(\frac{p_{t,1}}{p_{t-1,1}},\ldots,\frac{p_{t,N}}{p_{t-1,N}}\right)$ are the returns of the stocks
\item $f_t(\mathbf x)=-\log(\langle\mathbf x,\mathbf y_t\rangle)$ is the loss
\item $\mathcal E$ plays $\mathbf x_t=\mathbf x$ at each turn (Constant Rebalancing Portfolios). $\mathbf x^*=\arginf\limits_{x\in\Delta_{N-1}}\sum\limits_{t=1}^Tf_t(\mathbf x)$ is the Best Constant Rebalancing Portfolio
\end{itemize}

\onslide<2->{
\metroset{block=fill}
\begin{exampleblock}{Regret in Online Portfolio Optimization}
%The regret $R_T$ becomes:
$$R_T
%=\sup\limits_{\mathbf x\in\Delta_{N-1}}\sum\limits_{t=1}^T\log(\langle\mathbf x,\mathbf y_t \rangle)-\log(\langle\mathbf x_t,\mathbf y_t \rangle)
=\log(W_T(\mathbf x^*,\ldots,\mathbf x^*)/W_T(\mathbf x_1,\ldots,\mathbf x_T)),$$
where $W_T(\mathbf x_1,\ldots,\mathbf x_T)=\prod\limits_{t=1}^T\langle\mathbf x_t,\mathbf y_t\rangle$ is the wealth %obtained by playing the portfolio sequence $\mathbf x_1,\ldots,\mathbf x_T$
\end{exampleblock}
}

\onslide<3->{
\textbf{Limitations:} No transaction costs
}

\end{frame}

\begin{frame}{Adding Trading Costs}

% How much do we move during learning?
% How much do we pay in transaction costs?

\metroset{block=fill}
\begin{exampleblock}{Regret in Online Portfolio Optimization with Transaction Costs}

$$R^C_T
%=\sup\limits_{\mathbf x\in\Delta_{N-1}}\sum\limits_{t=1}^T\log(\langle\mathbf x,\mathbf y_t \rangle)-\log(\langle\mathbf x_t,\mathbf y_t \rangle)
=\underbrace{\log(W_T(\mathbf x^*,\ldots,\mathbf x^*)/W_T(\mathbf x_1,\ldots,\mathbf x_T))}_{R_T \text{: standard regret}}+\underbrace{\gamma\sum\limits_{t=1}^T||\mathbf x_t-\mathbf x_{t-1}||_1}_{C_T \text{: transaction costs}},$$
where $\gamma$ is the transaction rate
\end{exampleblock}

%This is derived naturally from proportional transaction costs:

% $$W_T(\mathbf x_1,\ldots,\mathbf x_T)=\prod\limits_{t=1}^T\alpha_t\langle\mathbf x_t,\mathbf y_t\rangle,$$ 
% where $\alpha_t\in[0,1]$ is the percentage of wealth left after incurring in some transaction costs. $\alpha_t$ is defined as: $\alpha_t:=1 - \gamma||\mathbf x_t - \mathbf x_{t-1}||_1$.

% Proportional Transaction costs: $\alpha_t:=1 - \gamma|| \alpha_t\mathbf x_t - \mathbf x'_{t-1}||_1$, $\mathbf x'_{t}=\frac{\mathbf x_t\otimes \mathbf y_t}{\langle\mathbf x_t,\mathbf y_t\rangle}$.

$\gamma$ is the proportional transaction rate for buying and selling stocks
\end{frame}

\section{State of the Art}

% \begin{frame}{State of the Art Algorithms}

% \begin{itemize}
%   \item Algorithms for Online Portfolio Optimization:
%   \begin{itemize}
%   \item Universal Portfolios (UP) [\cite{cover1996universal}]
%   \item Online Newton Step (ONS) [\cite{agarwal2006algorithms}]
%   \end{itemize}
%     \item Algorithms for Online Portfolio Optimization with transaction costs:
%   \begin{itemize}
%     \item Online Lazy Updates [\cite{das2013online}]
%   \end{itemize}
% \end{itemize}

% \end{frame}

\begin{frame}{State of the Art Algorithms}

\begin{itemize}
  \item \textbf{Online Portfolio Optimization}:
  \begin{itemize}
    \item \alert{Online Newton Step (ONS) [\cite{agarwal2006algorithms}]}
    \item \alert{Universal Portfolios (UP) [\cite{cover1996universal}]}
  \end{itemize}
  \item \textbf{Online Portfolio Optimization with Transaction Costs}:
  \begin{itemize}
    \item \alert{Online Lazy Updates (OLU) [\cite{das2013online}]}
  \end{itemize}
  \item \textbf{Heuristics}:
  \begin{itemize}
    \item Passive Aggressive Mean Reversion (PAMR) [\cite{li2012pamr}]
    \item Online Moving Average Reversion (OLMAR) [\cite{li2015moving}]
  \end{itemize}
\end{itemize}

\end{frame}

\section{Proposed Solution}

\begin{frame}{Online Gradient Descent}
\onslide<1->{
\renewcommand{\thealgorithm}{1}
\begin{algorithm}[H]
\small
    \caption{OGD in Online Portfolio Optimization with Transaction Costs}
    \label{alg:OGD_in_OPO}
    \begin{algorithmic}[1]
    \REQUIRE learning rate sequence $\{\eta_1, \ldots, \eta_T\}$  \nonumber
    \STATE Set $\mathbf{x}_1 \gets \frac{1}{N} \mathbf{1}$ \label{line:init_OGD}
    \FOR {$t \in \{ 1, \ldots, T \}$}
    \STATE $\mathbf{z}_{t+1} \gets \mathbf{x}_{t}+ \eta_t \frac{\mathbf{y}_t}{\langle \mathbf{y}_t,\mathbf{x}_t \rangle}$ \label{line:update_OGD}
    \STATE Select Portfolio $\mathbf x_{t+1}=\arginf\limits_{\mathbf x\in\Delta_{N-1}}||\mathbf z_t-\mathbf x||_2^2$ \label{line:line_projection_OGD}

    \STATE Observe $\mathbf{y}_{t+1}$ from the market \label{line:out_OGD}
    \STATE Get wealth $\log( \langle \mathbf{y}_{t+1},\mathbf{x}_{t+1} \rangle) - \gamma|| \mathbf{x}_{t+1} - \mathbf{x}_{t} ||_1$ \label{line:wealth}
    \ENDFOR
    \end{algorithmic}
\end{algorithm}
}
\vspace{-1.5em}
\begin{columns}[T] % align columns
\begin{column}{.50\textwidth}
\metroset{block=fill}
\begin{exampleblock}{Regret [\cite{zinkevich2003online}]}
\vspace{1.2em}
$$R_T\le\left( \frac{1}{K} + \frac{N K \epsilon_u^2}{\epsilon_l^2} \right) \sqrt{T}$$

\end{exampleblock}

\end{column}

\begin{column}{.50\textwidth}
\metroset{block=fill}
\begin{exampleblock}{Total Regret (this work)}

$$R_T^C\le \left[\hspace{-0.1em} \frac{1}{K}\hspace{-0.1em} +\hspace{-0.1em} \frac{NK\epsilon_u}{\epsilon_l}\hspace{-0.1em}\hspace{-0.1em}\left( \frac{\epsilon_u}{\epsilon_l} + 2 \gamma \right)\hspace{-0.2em} \right]\hspace{-0.3em} \sqrt{T}$$

\end{exampleblock}

\end{column}

\end{columns}



% \begin{figure}
% \centering
%  \scalebox{.55}{\begin{pspicture}(0,-2.35)(6.9,2.35)
\psellipse[linecolor=black, linewidth=0.04, dimen=outer](3.45,-0.4)(3.45,1.95)
\rput[bl](0.4,-0.55){$\mathcal D$}
\psdots[linecolor=black, dotsize=0.14078125](4.9,1.35)
\psline[linecolor=black, linewidth=0.04, arrowsize=0.05291667cm 2.0,arrowlength=1.4,arrowinset=0.0]{->}(4.9,1.35)(3.3,0.55)
\rput[bl](2.3,0.85){$\eta\nabla f_t(\mathbf x_t)$}
\rput[bl](4.7,1.45){$\mathbf x_t$}
\psline[linecolor=black, linewidth=0.04, arrowsize=0.05291667cm 2.0,arrowlength=1.4,arrowinset=0.0]{->}(4.9,1.35)(6.411111,2.15)
\rput[bl](6.6,2.15){$\mathbf z_t$}
\psdots[linecolor=black, dotsize=0.15374984](6.4,2.15)
\psline[linecolor=black, linewidth=0.04, arrowsize=0.05291667cm 2.0,arrowlength=1.4,arrowinset=0.0]{->}(6.3636365,2.0954545)(5.7272725,1.0045455)
\rput[bl](5.2363634,0.6227273){$\mathbf x_{t+1}$}
\end{pspicture}}
% \end{figure}

\end{frame}

% \begin{frame}{Theoretical Results of Online Gradient Descent}

% \begin{itemize}
% \item Standard Regret:
% $$R_T\le\left( \frac{1}{K} + \frac{N K \epsilon_u^2}{\epsilon_l^2} \right) \sqrt{T} $$
% \item Regret on the Costs: 
% $$\mathcolor{ultramarine}{C_T\le \frac{2 N K \gamma \epsilon_u}{\epsilon_l} \sqrt{T}}$$
% \item Total regret: 
% $$R_T^C\le \left[ \frac{1}{K} + \frac{NK\epsilon_u}{\epsilon_l}\left( \frac{\epsilon_u}{\epsilon_l} + 2 \gamma \right) \right] \sqrt{T}$$

% \end{itemize}

% \end{frame}

\begin{frame}{Theoretical and Computational Comparison}

Comparison of the theoretical guarantees and computational complexity among the selected algorithms

\begin{table}[ht!]\centering\small
\begin{tabular}{ |c||c|c|c|c| }
 %\hline
 %\multicolumn{5}{|c|}{Comparison} \\
 \hline
  & \textbf{OGD} & \textbf{UP} & \textbf{ONS} & \textbf{OLU}\\
 \hline
 \hline
 $R_T$ & $\mathcal O(\sqrt T)$  & $\mathcolor{ultramarine}{\mathcal O(\log T)} $ & $\mathcolor{ultramarine}{\mathcal O(\log T)}$ & $\mathcal O(\sqrt T)$ \\
 \hline
 $R^C_T$ & $\mathcolor{ultramarine}{\mathcal O(\sqrt T)}$ & - & - & $\mathcal O(T)$\footnote{Assuming $\gamma=const.$, if $\gamma\propto\sqrt T$ then $R_T^C=\mathcal O(\sqrt T)$}\\
 \hline
 Complexity & $\mathcolor{ultramarine}{\Theta(N)}$ & $\Theta (T^N)$ & $\Theta (N^2)$  & $\mathcolor{ultramarine}{\Theta (N)}$\\
 \hline
\end{tabular}
%\caption{Comparison between the theoretical guarantees of the selected algorithms .}\label{tab:comparison}
\end{table}


\end{frame}

% \begin{frame}{Online Mirror Descent and why OGD works}
% \end{frame}

\section{Experiments}

\begin{frame}{Experimental Setting}

\begin{table}[ht!]\centering\small
\begin{tabular}{ |c||c|c|c|c| }
 \hline
 \multicolumn{5}{|c|}{Datasets} \\
 \hline
 Name & Market &Year Span & Days & Assets\\
 \hline
 \textbf{NYSE(O)} & New York Stock Exchange  & 1962 - 1984  &5651&   36\\
 TSE & Toronto Stock Exchange & 1994 - 1998  & 1258   &88\\
 SP500 & Standard Poor's 500 & 1998 - 2003 & 1276&  25\\
 \hline
\end{tabular}
%\caption{Description of the main datasets used commonly in the Online Portfolio Optimization literature.}\label{tab:dataset}
\end{table}

\textbf{Performance Metrics}:

Sample for $100$ times, $5$ randomly drawn assets

\begin{itemize}
  \item Average Annual Percentage Yield (APY) defined as 
$$A(W_T)=W_T^{250/T} - 1$$
\item Normalized per round transaction costs $\frac{C_t(\mathfrak U)}{\gamma t}$
\end{itemize}
\end{frame}

\begin{frame}{Results on the NYSE(O) dataset}

\begin{figure}[t!]
    \input{figures/fig_w_decay_l1.tex}
\end{figure}

\end{frame}

\begin{frame}{Results on the NYSE(O) dataset}

\begin{figure}[t!]
    \pgfkeys{/pgf/number format/.cd,1000 sep={}}
\pgfplotsset{every tick label/.append style={font=\tiny}}
\begin{tikzpicture}

\colorlet{colorUP}{red}
\colorlet{colorOGD}{green!50.0!black}
\colorlet{colorONS}{blue}
\definecolor{colorOLU}{rgb}{0.75,0.75,0}
\begin{semilogyaxis}
    [
        height = 4.9cm,
        width = 6cm,
        legend cell align={left},
        legend style={at={(0.97,0.97)}, anchor=north east, draw=white!80.0!black},
        tick align=outside,
        tick pos=left,
        x grid style={white!69.01960784313725!black},
        xlabel={$t$},
        xmin=-100, xmax=5600.5,
        ymin=0.0005, ymax=0.07,
        xtick style={color=black},
        xtick = {1, 1500, 3000, 4500, 6000},
        y grid style={white!69.01960784313725!black},
        ylabel={$\frac{C_t(\mathfrak{U})}{t}$},
        yticklabel style={rotate=90},
        ytick = {0.001, 0.01},
        yticklabels = {$0.001$, $0.01$},
    ]


    %OGD
    \addplot [stack plots=false, semithick, colorOGD, mark = o]  table[x expr=\thisrowno{0}, y expr=\thisrowno{1}] {./figures/fig_costs/ogd_costs.txt};
    %\addlegendentry{OGD}
    \addplot[name path=A,fill=none, draw=none,stack plots=false,forget plot]  table[x expr=\thisrowno{0}, y expr=\thisrowno{2}]
    {./figures/fig_costs/ogd_costs.txt};
    \addplot[name path=B,fill=none, draw=none,stack plots=false,forget plot]  table[x expr=\thisrowno{0}, y expr=\thisrowno{3}]
    {./figures/fig_costs/ogd_costs.txt};
    \addplot[colorOGD, opacity=0.3,forget plot] fill between[of=A and B];
    
    %UP
    \addplot [stack plots=false, semithick, colorUP, mark = square]  table[x expr=\thisrowno{0}, y expr=\thisrowno{1}] {./figures/fig_costs/up_costs.txt};
    %\addlegendentry{U$_c$P}
    \addplot[name path=A,fill=none, draw=none,stack plots=false,forget plot]  table[x expr=\thisrowno{0}, y expr=\thisrowno{2}] {./figures/fig_costs/up_costs.txt};
    \addplot[name path=B,fill=none, draw=none,stack plots=false,forget plot]  table[x expr=\thisrowno{0}, y expr=\thisrowno{3}] {./figures/fig_costs/up_costs.txt};
    \addplot[colorUP, opacity=0.3,forget plot] fill between[of=A and B];

    %OLU
    \addplot [stack plots=false, semithick, colorOLU, mark = diamond]  table[x expr=\thisrowno{0}, y expr=\thisrowno{1}] {./figures/fig_costs/olu_costs.txt};
    %\addlegendentry{OLU}
    \addplot[name path=A,fill=none, draw=none,stack plots=false,forget plot]  table[x expr=\thisrowno{0}, y expr=\thisrowno{2}] {./figures/fig_costs/olu_costs.txt};
    \addplot[name path=B,fill=none, draw=none,stack plots=false,forget plot]  table[x expr=\thisrowno{0}, y expr=\thisrowno{3}] {./figures/fig_costs/olu_costs.txt};
    \addplot[colorOLU, opacity=0.3,forget plot] fill between[of=A and B];
    
    %ONS
    \addplot [stack plots=false, semithick, colorONS, mark = triangle]  table[x expr=\thisrowno{0}, y expr=\thisrowno{1}] {./figures/fig_costs/ons_costs.txt};
    %\addlegendentry{ONS}
    \addplot[name path=A,fill=none, draw=none,stack plots=false,forget plot]  table[x expr=\thisrowno{0}, y expr=\thisrowno{2}] {./figures/fig_costs/ons_costs.txt};
    \addplot[name path=B,fill=none, draw=none,stack plots=false,forget plot]  table[x expr=\thisrowno{0}, y expr=\thisrowno{3}] {./figures/fig_costs/ons_costs.txt};
    \addplot[colorONS, opacity=0.3, forget plot] fill between[of=A and B];
\end{semilogyaxis} \end{tikzpicture}

\end{figure}

\end{frame}

\section{Conclusions}

\begin{frame}{Conclusions and Future Work}

\textbf{Contributions:}
\begin{itemize}
\item OGD in online portfolio optimization
\item Analysis of the total regret of OGD
\item Experimental campaign on real data
\end{itemize}

\textbf{Future Works:}

\begin{itemize}
  \item Develop cost aware algorithms (work in progress)
  \item Generalize the regret analysis to other online learning algorithms
\end{itemize}

\end{frame}

\begin{frame}[allowframebreaks]
  \bibliographystyle{apalike}
  %\bibliographystyle{abbrv}
  \bibliography{bibl_thesis.bib}
  

\end{frame}

\section{Appendix}

\appendix

\begin{frame}{Results on the SP500 dataset (proportional transaction costs)}

\begin{figure}[t!]
    \input{figures/fig_w_decay_true.tex}
\end{figure}
\end{frame}

\begin{frame}{Results on the TSE dataset}

\begin{figure}[t!]
    \input{figures/fig_w_decay_tse.tex}
\end{figure}

\end{frame}

\begin{frame}{Results on the SP500 dataset}

\begin{figure}[t!]
    \pgfplotsset{every tick label/.append style={font=\tiny}}
\colorlet{colorUP}{red}
\colorlet{colorOGD}{green!50.0!black}
\colorlet{colorONS}{blue}
\definecolor{colorOLU}{rgb}{0.75,0.75,0}
\begin{tikzpicture}

    \begin{semilogxaxis}
    [
        height = 8cm,
        width = 9cm,
        legend cell align={left},
        legend style={at={(0.03,0.03)}, anchor=south west, draw=white!80.0!black},
        tick align=outside,
        tick pos=left,
        x grid style={white!69.01960784313725!black},
        xlabel={$\gamma$},
        xmin=0.00009, xmax=0.045,
        ymin=-0.01, ymax=0.2,
        xtick style={color=black},
        xtick ={0.0001,0.001,0.006,0.04},
        xticklabels={$0$, $0.001$, $0.006$, $0.04$},
        y grid style={white!69.01960784313725!black},
        ylabel={$A(W_T)$},
        ytick style={color=black},
        yticklabels={$-0.15$, $-0.1$, $-0.05$, $0.0$, $0.05$, $0.1$, $0.15$, $0.2$},
        ytick = {-0.15,-0.1,-0.05,0,0.05,0.1,0.15,0.2},
        ymin=-0.15, ymax=0.2,
        yticklabel style={rotate=90},
    ]

    %OGD
    \addplot [stack plots=false, semithick, colorOGD, mark = o]  table[x expr=\thisrowno{0}, y expr=\thisrowno{1}] {./figures/fig_wealth_decay_sp/ogd_costs_l1.txt};
    \addlegendentry{OGD}
    \addplot[name path=A,fill=none, draw=none,stack plots=false,forget plot]  table[x expr=\thisrowno{0}, y expr=\thisrowno{2}]
    {./figures/fig_wealth_decay_sp/ogd_costs_l1.txt};
    \addplot[name path=B,fill=none, draw=none,stack plots=false,forget plot]  table[x expr=\thisrowno{0}, y expr=\thisrowno{3}]
    {./figures/fig_wealth_decay_sp/ogd_costs_l1.txt};
    \addplot[colorOGD, opacity=0.3,forget plot] fill between[of=A and B];
    
    %UP
    \addplot [stack plots=false, semithick, colorUP, mark = square]  table[x expr=\thisrowno{0}, y expr=\thisrowno{1}] {./figures/fig_wealth_decay_sp/up_costs_l1.txt};
    \addlegendentry{UP}
    \addplot[name path=A,fill=none, draw=none,stack plots=false,forget plot]  table[x expr=\thisrowno{0}, y expr=\thisrowno{2}] {./figures/fig_wealth_decay_sp/up_costs_l1.txt};
    \addplot[name path=B,fill=none, draw=none,stack plots=false,forget plot]  table[x expr=\thisrowno{0}, y expr=\thisrowno{3}] {./figures/fig_wealth_decay_sp/up_costs_l1.txt};
    \addplot[colorUP, opacity=0.3,forget plot] fill between[of=A and B];
    
    %OLU
    \addplot [stack plots=false, semithick, colorOLU, mark = diamond]  table[x expr=\thisrowno{0}, y expr=\thisrowno{1}] {./figures/fig_wealth_decay_sp/olu_costs_l1.txt};
    \addlegendentry{OLU}
    \addplot[name path=A,fill=none, draw=none,stack plots=false,forget plot]  table[x expr=\thisrowno{0}, y expr=\thisrowno{2}] {./figures/fig_wealth_decay_sp/olu_costs_l1.txt};
    \addplot[name path=B,fill=none, draw=none,stack plots=false,forget plot]  table[x expr=\thisrowno{0}, y expr=\thisrowno{3}] {./figures/fig_wealth_decay_sp/olu_costs_l1.txt};
    \addplot[colorOLU, opacity=0.3,forget plot] fill between[of=A and B];
    
    %ONS
    \addplot [stack plots=false, semithick, colorONS, mark = triangle]  table[x expr=\thisrowno{0}, y expr=\thisrowno{1}] {./figures/fig_wealth_decay_sp/ons_costs_l1.txt};
    \addlegendentry{ONS}
    \addplot[name path=A,fill=none, draw=none,stack plots=false,forget plot]  table[x expr=\thisrowno{0}, y expr=\thisrowno{2}] {./figures/fig_wealth_decay_sp/ons_costs_l1.txt};
    \addplot[name path=B,fill=none, draw=none,stack plots=false,forget plot]  table[x expr=\thisrowno{0}, y expr=\thisrowno{3}] {./figures/fig_wealth_decay_sp/ons_costs_l1.txt};
    \addplot[colorONS, opacity=0.3, forget plot] fill between[of=A and B];
    
    \end{semilogxaxis}

\end{tikzpicture}
\end{figure}

\end{frame}

\begin{frame}{Comparison of Algorithms (single run of NYSE(O))}

\begin{figure}[t!]
    \begin{tikzpicture}

\colorlet{colorUP}{red}
\colorlet{colorOGD}{green!50.0!black}
\colorlet{colorONS}{blue}
\definecolor{colorOLU}{rgb}{0.75,0.75,0}
\pgfkeys{/pgf/number format/.cd,1000 sep={}}
\pgfplotsset{every tick label/.append style={font=\tiny}}
\begin{axis}[
height = 4. cm,
width = 5cm,
legend cell align={left},
legend style={at={(0.03,0.97)}, anchor=north west, draw=white!80.0!black},
tick align=outside,
tick pos=left,
x grid style={white!69.01960784313725!black},
xlabel={$t$},
xmin=0.5, xmax=5700.5,
ymin=0, ymax=40,
xtick style={color=black},
xtick = {1, 1000, 2000, ..., 5000, 6000},
y grid style={white!69.01960784313725!black},
ylabel={${W}^C_t(\mathfrak{U})$},
ytick style={color=black},
yticklabel style={rotate=90},
]
\addplot [thick, colorOGD, mark = o,mark repeat=20,mark phase=7]
table {./figures/fig_11/small/data_ogd.txt};
\addlegendentry{\small OGD}
\addplot [thick, colorUP, mark = square,mark repeat=20,mark phase=7]
table {./figures/fig_11/small/data_up.txt};
\addlegendentry{\small U$_C$P}
\addplot [semithick, colorOLU, mark = diamond,mark repeat=20,mark phase=7]
table {./figures/fig_11/small/data_olu.txt};
\addlegendentry{\small OLU}
\addplot [semithick, colorONS, mark = triangle,mark repeat=20,mark phase=7]
table {./figures/fig_11/small/data_ons.txt};
\addlegendentry{\small ONS}
%\addplot [semithick, magenta, mark = +,mark repeat=20,mark phase=7]
%table {./figures/fig_11/small/data_bah.txt};
%\addlegendentry{B\&H}
\end{axis}

\end{tikzpicture}

\end{figure}

\end{frame}


\end{document}
